% Source: https://artofproblemsolving.com/community/c3h2397502_happy_new_year_2021
% Stable identity: aops-topic-2397502-q04
\textbf{A particular solution.} $P(x)=x$ works: $x(x+1)=(x+2)x-x$.

\textbf{The difference.} Let $Q(x)=P(x)-x$. Subtracting the identity for $x$ from the given one,
\[xQ(x+1)=(x+2)Q(x)\qquad\text{for all }x.\]
Put $x=0$: $0=2Q(0)$, so $Q(0)=0$. Put $x=-1$: $-Q(0)=Q(-1)$, so $Q(-1)=0$. By the factor theorem $Q(x)=x(x+1)R(x)$ for a polynomial $R$. Substituting,
\[x(x+1)(x+2)R(x+1)=(x+2)x(x+1)R(x),\]
so $R(x+1)=R(x)$ for every $x\notin\{0,-1,-2\}$, and hence (a polynomial with infinitely many roots is zero) for all $x$. Then $R(n)=R(0)$ for every positive integer $n$, so the polynomial $R(x)-R(0)$ has infinitely many roots and $R$ is a constant $c$.

\textbf{Answer.} $P(x)=x+cx(x+1)=\boxed{cx^2+(c+1)x}$ for any real constant $c$.

\textbf{Check.} $xP(x+1)=x(x+1)+cx(x+1)(x+2)$ and $(x+2)P(x)-x=x^2+2x+cx(x+1)(x+2)-x=x(x+1)+cx(x+1)(x+2)$.
